%% HAND-WRITTEN table fragment. Sizes assume eight stored gain codes per
%% channel; the per-code reference-DAC residual table is excluded on the
%% argument in note_calibration_residual (it buys 0.07 mV rms for 4 KiB).
\begin{table}[htbp]
  \centering
  \caption[{Per-chip calibration constants.}]{Per-chip calibration constants,
  their origin and their refresh rate. \emph{Test} constants need external
  equipment and are written once; \emph{power-up} constants are re-measured
  on-die at every start and overwrite the stored copy. At eight stored gain codes
  per channel --- the full set a channel uses over its concentration range ---
  the whole block is under 128\,bytes per channel, so
  it is staged from flash into a channel hart's private memory with its image
  rather than read back during a measurement.}
  \label{tab:calibration-constants}
  \small
  \begin{tabular}{llll}
    \toprule
    Constant & Scope & Set at & Refreshed \\
    \midrule
    Bias trim code & chip & test & --- \\
    Reference DAC offset and step & channel & test & --- \\
    Reference droop parameter & channel & test & --- \\
    Converter step size & channel & test & --- \\
    Working-electrode switch resistance & channel & test & --- \\
    Transimpedance gain, per gain code & channel & test & scaled at power-up \\
    \addlinespace
    Converter zero code & channel & test & power-up, and between steps \\
    A1 offset and tracking slope & channel & test & power-up \\
    Current zero, per gain code & channel & test & power-up \\
    Internal gain ratio, per gain code & channel & test & power-up \\
    \bottomrule
  \end{tabular}
\end{table}
